\documentclass[11pt]{article}
\usepackage[utf8]{inputenc}
\usepackage[T1]{fontenc}
\usepackage{amsmath,amssymb}
\usepackage[brazil]{babel}

\title{LeetCode 504 --- Base 7}
\author{}
\date{}

\begin{document}
\maketitle

\section{Representação posicional}

Para $n \ge 0$, existe uma sequência única de dígitos $d_i \in \{0,1,2,3,4,5,6\}$ tal que
\[
  n = \sum_{i=0}^{k} d_i \, 7^{i},
\]
com $d_k \neq 0$ quando $n > 0$. Caso especial: $n = 0 \Leftrightarrow d_0 = 0$ (string \texttt{"0"}).

O índice $i = 0$ é o dígito das unidades ($7^0$), ou seja, o caractere mais à \textbf{direita} na string.

\section{Extração dos dígitos (restos)}

\[
\begin{aligned}
  r_0 &= n \bmod 7, & q_0 &= \left\lfloor \frac{n}{7} \right\rfloor, \\
  r_1 &= q_0 \bmod 7, & q_1 &= \left\lfloor \frac{q_0}{7} \right\rfloor, \\
  r_2 &= q_1 \bmod 7, & q_2 &= \left\lfloor \frac{q_1}{7} \right\rfloor, \\
  &\vdots & &
\end{aligned}
\]

Defina $d_i = r_i$. Pare quando $q_j = 0$. A string é $d_k d_{k-1} \cdots d_0$ ($d_k$ à esquerda).

No loop, $d_0$ sai primeiro; concatenar na ordem do loop deixa a string \textbf{invertida} em relação ao resultado final.

\section{Exemplo $n = 100$}

\[
\begin{aligned}
  100 &= 14 \cdot 7 + 2 &&\Rightarrow d_0 = 2, \\
  14 &= 2 \cdot 7 + 0 &&\Rightarrow d_1 = 0, \\
  2 &= 0 \cdot 7 + 2 &&\Rightarrow d_2 = 2.
\end{aligned}
\]

Verificação:
\[
  100 = 2 \cdot 7^{2} + 0 \cdot 7^{1} + 2 \cdot 7^{0} = 98 + 0 + 2.
\]
String: \texttt{"202"}.

\section{Inteiros negativos}

\[
  n < 0 \quad \Rightarrow \quad \text{resposta} = \texttt{"-"} \mathbin{\|} \text{base}_7(|n|).
\]
Exemplo: $-7 \mapsto \texttt{"-10"}$.

\section{Contagem decimal $1$--$10$}

\begin{center}
\begin{tabular}{r|l}
  Decimal & Base 7 \\ \hline
  1  & 1 \\
  2  & 2 \\
  3  & 3 \\
  4  & 4 \\
  5  & 5 \\
  6  & 6 \\
  7  & 10 \\
  8  & 11 \\
  9  & 12 \\
  10 & 13 \\
\end{tabular}
\end{center}

\end{document}
